\section*{Introduction}

The size of a population is a key factor mediating the mode and tempo of evolution,
particularly as it relates to the strength of stochastic versus deterministic
evolutionary processes \citep{kimura_mutation_1963,gillespie_genetic_2000}.
Population genetics theory predicts that selection is less efficacious in small populations
\citep{charlesworth_effective_2009,gravel_when_2016}, resulting in reduced genetic
diversity from genetic drift and an increase in genetic load due to the accumulation
of deleterious mutations \citep{lande_risk_1994,lynch_mutation_1995,dussex_purging_2023}.
These dynamics have been extensively observed in natural populations, e.g., the
increase in genetic load in humans as a function of their geographic distance to more
diverse groups in sub-Saharan Africa \citep{henn_distance_2016} and the increase in
genetic load and inbreeding depression in small populations of conservation concern
\citep{lande_genetics_1988,westemeier_prairiechicken_1998,frankham_genetics_2005,
kardos_inbreeding_2023}, empirically supporting the theoretical expectations.
\vspace{0.2cm}

In contrast, evolution in large populations is thought to be primarily mediated by
deterministic processes due to an increase in the efficacy of selection
\citep{charlesworth_effective_2009,gravel_when_2016,lanfear_population_2014}.
The dynamics of selection in large populations are expected to produce patterns
such as clonal-interference effects driven by soft sweeps and recurrent mutation
\citep{park_clonal_2007,weissman_rate_2014}, as well as the loss of genetic
diversity through the hitchhiking effect \citep{smith_hitch-hiking_1974}.
In sufficiently large populations, this process is predicted to be a primary driver
of stochastic diversity loss \citep{gillespie_genetic_2000,gillespie_role_1999}.
In stark contrast to the well-studied dynamics of small populations, empirical work
on these mechanisms has been limited, with a few key exceptions
\citep{dey_molecular_2013,ag1000g,teterina_pervasive_2025,anderson_population_2017,
braso-vives_highly_2026}.
Empirical studies on this subject have mainly focused on eukaryotic species with
small-to-moderate historical population sizes \citep{cutter_molecular_2013,
leffler_revisiting_2012,buffalo_quantifying_2021}, limiting our capacity to observe
these evolutionary dynamics at demographic extremes.
This gap highlights the need for empirical studies in systems with historically large
population sizes and elevated levels of genetic polymorphism.
\vspace{0.2cm}

The Pacific acorn barnacle (\balgla{}; Balanidae) presents an exceptional model for studying the
dynamics and mechanisms of evolution in large populations. This species is widespread
across the Pacific intertidal of North America, from Baja California to Alaska,
where it can reach densities of tens of thousands of individuals per square meter
\citep{menge_recruitment_2000}. Like most acorn barnacles, \bgla{} is a sessile, simultaneous hermaphrodite that
reproduces through cross-fertilization, mating with its neighbors via pseudo-copulation;
even isolated individuals, beyond the reach of any mate, can produce outcrossed broods
by capturing water-borne sperm \citep{barazandeh_spermcast_2014}.
Individuals produce several large broods per year \citep{hines_reproduction_1978},
releasing nauplius larvae that feed and develop in the plankton for a period of weeks
before settling \citep{brown_larvae_1985,pfeiffer-hoyt_modeling_2005}.
This combination of high fecundity, outcrossing, and long-lived planktonic larvae
capable of long-distance dispersal, together with these enormous census sizes, is
consistent with very large effective population sizes, making \bgla{} an ideal system
for testing theoretical predictions about selection and drift in large populations.
\vspace{0.2cm}

This species has been the focus of extensive ecological and genetic research. Prior work
has linked environmental factors to variation in development, recruitment, dispersal,
and reproduction \citep{berger_spatial_2006,berger_reproduction_2009,
barshis_coastal_2011,pfeiffer-hoyt_modeling_2005,schubart_recruitment_1995}, and
substantial effort has characterized the spatial population structure of \bgla{}
along the Pacific coast, particularly a steep genetic cline on the coast of
California that separates northern and southern populations
\citep{sotka_strong_2004,wares_diversification_2005,wares_maintenance_2019,
wares_diversity_2021}.
Yet, \bgla{} remains genomically understudied.
Previous research has been limited to a handful of nuclear and mitochondrial markers,
with the most extensive study employing reduced representation sequencing in the
absence of a reference genome \citep{wares_diversity_2021}.
Moreover, while two reference assemblies for \bgla{} exist in public databases, both are
highly fragmented and likely incomplete, which limits their use when studying genome-wide
patterns of diversity and selection in this system.
\vspace{0.2cm}

Most publicly available barnacle assemblies suffer from similar limitations, presumably
due to the difficulty of assembling highly polymorphic genomes from short reads
\citep{vinson_assembly_2005,kajitani_efficient_2014}.
Despite this limitation, a handful of chromosome-level assemblies have been recently
generated for several barnacle species \citep{han_new_2024,bishop_genome_2025,
bernot_chromosome-level_2022,carlson_reference_2025,blaxter_genome_2023}.
Yet, these genomes largely belong to evolutionarily or geographically distant species.
While useful for the long-term comparisons we perform here, their evolutionary distance
limits the utility of these resources for studies in \bgla{}.
This drives the need for high-quality assemblies of both evolutionarily
closely related species, which could serve as outgroups for recent evolutionary
comparisons, and co-localized species, which enable the study of ecological factors
maintaining diversity in this system.
\vspace{0.2cm}

Work in other barnacle species has revealed substantial levels of genetic diversity
\citep{alm_rosenblad_genomic_2021}, in accordance with the expectation of very large
effective populations, and complex patterns of selection at remarkably small spatial
scales.
The best-characterized example is the balanced polymorphism at the mannose-6-phosphate
isomerase (\textit{Mpi}) locus in the northern acorn barnacle \textit{Semibalanus
balanoides}, which segregates across intertidal environments
\citep{schmidt_environmental_2000,nunez_footprints_2020}.
Building on this single-locus work, whole-genome pooled resequencing across intertidal
microhabitats has shown that spatially varying selection extends beyond \textit{Mpi},
leaving genomic signatures of adaptation to environmental heterogeneity at loci
throughout the \textit{S. balanoides} genome \citep{nunez_tides_2021}.
These studies demonstrate that selection can structure barnacle genetic variation over
just meters of tidal height; however, such genome-wide inference has so far relied on
a fragmented draft assembly in a single species, and the landscape of diversity and
selection remains uncharacterized across most of the barnacle tree, including
\bgla{}.
\vspace{0.2cm}

In this work we start addressing the need for empirical systems to study evolution in
large populations by generating a high-quality, chromosome-level assembly for the Pacific
acorn barnacle.
We additionally generate a draft assembly for its congener, the wrinkled barnacle
\textit{Balanus crenatus}, providing a closely related outgroup for divergence-based analyses.
We use these assemblies to perform whole-genome comparative analyses, revealing conservation
in genome architecture across 165 million years of barnacle evolution.
We further conduct the first genome-wide assessment of genetic variation in this system,
uncovering extreme levels ($>5\%$) of nucleotide diversity, with over 20 million SNPs
observed in just six individuals.
Consistent with theoretical expectations, we observe reduced diversity in protein-coding
sequences, reflecting strong background selection and functional
constraints.
McDonald-Kreitman tests further reveal that the majority of non-synonymous substitutions
between species were driven by positive selection, as expected when selection is highly
efficacious in large populations.
In addition to providing high-quality genomic resources for a genomically understudied
group, our assemblies lay the groundwork for future work in \bgla{} and other barnacles as
models for the study of evolution in large, hyperpolymorphic populations.
\vspace{0.2cm}

